\documentclass[a4paper]{article}
\usepackage{graphicx}
\usepackage{amsmath,amssymb}
% web referencing
\usepackage[colorlinks=true, linkcolor=blue, urlcolor=blue]{hyperref}
\newcommand{\seehere}[1]{\href{#1}{see here}
\usepackage{minted}
\usepackage{multirow}
\usepackage{float}
\usepackage{todonotes}
\usetikzlibrary{graphs,graphdrawing}
\usegdlibrary{layered}
\usepackage{forest}
\usepackage{xstring}
\input{/home/donkarlo/Dropbox/repo/nd_ai_project/src/nd_ai/knoweldge_representation/language/formal/computer/programming/tex/latex/code_clips/external_dot.tex}
\title{}
\date{}

\begin{document}
    \maketitle
    A cognitive architecture is both a theory about the structure of the human mind and a computational instantiation of such a theory used in the fields of artificial intelligence (AI) and computational cognitive science. \url{https://en.wikipedia.org/wiki/Cognitive_architecture}
    
    For diffrent type \
    % \bibliographystyle{plain}
    % \bibliography{/home/donkarlo/Dropbox/repo/shared_research/bibliography/bibliography.bib}
\end{document}